% =====================================================================
%  BÖLÜM 4: Beklenti ve Momentler
%  Kaynak: Feller; Billingsley; Klenke; Shiryaev
% =====================================================================
\chapter{Beklenti ve Momentler}
\label{chp:beklenti}

\bolumalinti{%
  Beklenti, belirsizliğin ağırlık merkezidir.%
}{Patrick Billingsley, \textit{Probability and Measure}}

% ---- Kazanımlar (TYMM) ----
\begin{kazanimlar}
Bu bölümün sonunda öğrenci:
\begin{itemize}[leftmargin=1.8em]
  \item Beklentiyi Lebesgue integrali olarak tanımlayabilir ve
        bu tanımın kesikli/sürekli formüllere nasıl indirgeneceğini gösterebilir.
  \item Beklentinin doğrusallık, monotonluk ve önemli eşitsizliklerini
        (Jensen, Markov, Chebyshev) ispatlayabilir ve uygulayabilir.
  \item Varyans, kovaryans ve korelasyon kavramlarını hesaplayabilir;
        bağımsızlıkla ilişkilerini açıklayabilir.
  \item Moment üreten fonksiyon (MÜF) ve karakteristik fonksiyonun (KF)
        tanımını, yakınsama bölgesini ve momentlerle ilişkisini açıklayabilir.
  \item Koşullu beklentiyi $L^2$ projeksiyonu olarak yorumlayabilir,
        temel özelliklerini ve yineleme (kule) özelliğini ispatlayabilir.
\end{itemize}
\end{kazanimlar}

% ---- Motivasyon ----
\begin{motivasyon}
Bir rasgele değişkenin dağılımı tam bilgiyi içerir; ama çoğu zaman bu kadar bilgiye
ihtiyaç duymayız. Dağılımı birkaç sayıyla özetlemek istiyoruz: "ortalama değer nedir?"
ve "ne kadar saçılım var?" Beklenti ve varyans bu iki soruyu yanıtlar.

Daha incelikli özetler için momentler, çarpıklık (skewness) ve basıklık (kurtosis)
devreye girer. Tüm momentleri birden kodlayan moment üreten fonksiyon ve
karakteristik fonksiyon ise dağılım teorisinin (Böl.~\ref{chp:dagilim_aileleri})
ve Merkezi Limit Teoremi'nin (Böl.~\ref{chp:merkezi_limit}) temel araçlarıdır.

\medskip
\textbf{Köprü.} Bölüm~\ref{chp:olcum_teorisi}'deki Lebesgue integrali (MYT, Fatou, DYT)
burada olasılık dili ile yeniden okunur.
Bölüm~\ref{chp:rasgele_degiskenler}'deki PDF ve PMF kavramları, beklenti hesabını
somutlaştırır. Bölüm~\ref{chp:buyuk_sayilar}'da Büyük Sayılar Kanunu bu bölümün
sonuçlarına dayanacaktır.
\end{motivasyon}

% =====================================================================
\section{Beklenti: Lebesgue İntegrali Olarak}
\label{sec:beklenti_tanim}
% =====================================================================

\begin{kesif}
Altı yüzlü zarın beklenen değerini klasik yöntemle hesaplayalım:
$\frac{1}{6}(1+2+3+4+5+6) = 3.5$.
Bu hesap $\sum_x x \cdot P(X=x)$ formülünü izliyor.
Peki ya sürekli dağılımlar? Ya da karma dağılımlar?
Tüm bu durumları kapsayan tek bir tanım mümkün müdür?
\end{kesif}

\begin{tanimgerekce}
Bölüm~\ref{chp:olcum_teorisi}'de basit fonksiyonlarla başlayıp
non-negatif ve ardından genel fonksiyonlara geçen üç aşamalı Lebesgue
integral inşasını hatırlayın. Beklenti, olasılık ölçüsüne göre bu
integralin ta kendisidir — hem kesikli hem sürekli durumu tek çatı altında toplar.
\end{tanimgerekce}

\begin{tanim}[Beklenti]
\label{def:beklenti}
$(\Omega, \mathcal{F}, P)$ üzerinde $X$ rasgele değişken olsun.
$\int_\Omega |X| \, dP < \infty$ ise $X$'in \textbf{beklentisi} (expected value)
\[
  E[X] = \int_\Omega X(\omega) \, dP(\omega)
\]
olarak tanımlanır; $E[X]$ Bölüm~\ref{chp:olcum_teorisi}'deki Lebesgue integralidir.
\end{tanim}

\begin{teorem}[Beklentinin Somut Formülleri]
\label{thm:beklenti_formul}
$X$ rasgele değişken ve $E[|X|] < \infty$ olsun.
\begin{enumerate}[(i)]
  \item \textbf{Kesikli ($S = \{x_1, x_2, \ldots\}$):}
        $E[X] = \sum_k x_k \, P(X = x_k)$.
  \item \textbf{Mutlak sürekli (PDF $f_X$):}
        $E[X] = \int_{-\infty}^{+\infty} x \, f_X(x) \, dx$.
  \item \textbf{Bilinç değişimi (transfer teoremi):}
        $g : \mathbb{R} \to \mathbb{R}$ ölçülebilir ve $E[|g(X)|] < \infty$ ise
        \[
          E[g(X)] = \int_\Omega g(X(\omega)) \, dP(\omega) = \int_{\mathbb{R}} g(x) \, dP_X(x),
        \]
        yani sürekli durumda $E[g(X)] = \int g(x) f_X(x) \, dx$.
\end{enumerate}
\end{teorem}

\begin{proof}
\textit{(i):} $X = \sum_k x_k \mathbf{1}_{\{X = x_k\}}$ ayrışımı; Lebesgue integralinin doğrusallığı.\\
\textit{(ii):} $P_X(B) = \int_B f_X \, d\lambda$ (mutlak süreklilik);
$\int x \, dP_X(x) = \int x f_X(x) \, d\lambda(x)$.\\
\textit{(iii):} $h = g \circ X$ bileşkesi ölçülebilir; ölçü değişim teoremi (push-forward):
$\int_\Omega h \, dP = \int_{\mathbb{R}} g \, dP_X$.
\end{proof}

\begin{ornek}[Geometrik Dağılım]
$X \sim \mathrm{Geom}(p)$: $P(X = k) = (1-p)^{k-1} p$, $k \geq 1$.
\[
  E[X] = \sum_{k=1}^\infty k(1-p)^{k-1} p = p \cdot \frac{d}{dp}\!\left[-\sum_{k=1}^\infty (1-p)^k\right]
  = p \cdot \frac{d}{dp}\!\left[-\frac{1-p}{p}\right] = \frac{1}{p}.
\]
\end{ornek}

\begin{ornek}[Üstel Dağılım]
$X \sim \mathrm{Exp}(\lambda)$: $f_X(x) = \lambda e^{-\lambda x}$, $x > 0$.
\[
  E[X] = \int_0^\infty x \lambda e^{-\lambda x} \, dx
  = \left[-x e^{-\lambda x}\right]_0^\infty + \int_0^\infty e^{-\lambda x} \, dx
  = 0 + \frac{1}{\lambda} = \frac{1}{\lambda}.
\]
\end{ornek}

\begin{hataanalizi}
"Her rasgele değişkenin beklentisi vardır" — yanlış.
Cauchy dağılımı ($f(x) = \frac{1}{\pi(1+x^2)}$) için $\int_{-\infty}^{+\infty} |x| f(x) \, dx = \infty$;
beklenti tanımsızdır. St.\ Petersburg paradoksu (bir madeni para oyununda beklenti sonsuz
ama oyunu oynamak istediğiniz ücret sonludur) da bu olgunun pratik yansımasıdır.
\end{hataanalizi}

\begin{onerme}[Kuyruk Formülü]
\label{prop:kuyruk_beklenti}
$X \geq 0$ rasgele değişken için
\[
  E[X] = \int_0^\infty P(X > t) \, dt.
\]
Daha genel: $E[X^r] = r \int_0^\infty t^{r-1} P(X > t) \, dt$ ($r > 0$).
\end{onerme}

\begin{proof}
Fubini:
\[
  \int_0^\infty P(X > t) \, dt = \int_0^\infty \int_\Omega \mathbf{1}_{X(\omega) > t} \, dP(\omega) \, dt
  = \int_\Omega \int_0^{X(\omega)} dt \, dP(\omega) = \int_\Omega X \, dP = E[X].
\]
\end{proof}

\begin{mikroozet}
$E[X] = \int X \, dP$; kesikli: $\sum x_k p_k$; sürekli: $\int x f(x) \, dx$.
Her rasgele değişkenin beklentisi olmayabilir ($E[|X|] < \infty$ gereklidir).
\end{mikroozet}

% =====================================================================
\section{Beklentinin Temel Özellikleri}
\label{sec:beklenti_ozellik}
% =====================================================================

\begin{teorem}[Beklentinin Temel Özellikleri]
\label{thm:beklenti_ozellik}
$X, Y$ rasgele değişkenler, $E[|X|], E[|Y|] < \infty$, $a, b \in \mathbb{R}$:
\begin{enumerate}[(i)]
  \item \textbf{Doğrusallık:} $E[aX + bY] = aE[X] + bE[Y]$.
  \item \textbf{Monotonluk:} $X \leq Y$ n.k.\ ise $E[X] \leq E[Y]$.
  \item \textbf{$|E[X]| \leq E[|X|]$} (mutlak değer eşitsizliği).
  \item \textbf{Sabit:} $E[c] = c$.
  \item \textbf{Bağımsızlık:} $X, Y$ bağımsız ise $E[XY] = E[X] \cdot E[Y]$
        (her iki beklenti de mevcutsa).
\end{enumerate}
\end{teorem}

\begin{proof}
\textit{(i)}: Lebesgue integralinin doğrusallığı.\\
\textit{(ii)}: $X \leq Y$ n.k.\ $\Rightarrow$ $Y - X \geq 0$ n.k.\ $\Rightarrow$ $E[Y-X] \geq 0$.\\
\textit{(iii)}: $-|X| \leq X \leq |X|$; monotonluktan $|E[X]| \leq E[|X|]$.\\
\textit{(v)}: Sürekli durumda Fubini:
$E[XY] = \iint xy \, f_X(x) f_Y(y) \, dx \, dy = \left(\int x f_X(x) \, dx\right)\!\left(\int y f_Y(y) \, dy\right)
= E[X] \cdot E[Y]$.
\end{proof}

\begin{teorem}[Markov Eşitsizliği]
\label{thm:markov}
$X \geq 0$ rasgele değişken ve $a > 0$ için
\[
  P(X \geq a) \leq \frac{E[X]}{a}.
\]
\end{teorem}

\begin{proof}
$X \geq X \cdot \mathbf{1}_{X \geq a} \geq a \cdot \mathbf{1}_{X \geq a}$.
Beklenti alınca: $E[X] \geq a \cdot P(X \geq a)$. Düzenlenir.
\end{proof}

\begin{teorem}[Chebyshev Eşitsizliği]
\label{thm:chebyshev}
$X$ rasgele değişken, $E[X] = \mu$, $\mathrm{Var}(X) = \sigma^2 < \infty$. Her $k > 0$ için
\[
  P(|X - \mu| \geq k\sigma) \leq \frac{1}{k^2}.
\]
Eşdeğer form: $P(|X - \mu| \geq \varepsilon) \leq \frac{\sigma^2}{\varepsilon^2}$.
\end{teorem}

\begin{proof}
$(X-\mu)^2 \geq 0$ rasgele değişkenine $a = \varepsilon^2$ ile Markov:
$P((X-\mu)^2 \geq \varepsilon^2) \leq E[(X-\mu)^2]/\varepsilon^2 = \sigma^2/\varepsilon^2$.
$P((X-\mu)^2 \geq \varepsilon^2) = P(|X-\mu| \geq \varepsilon)$.
\end{proof}

\begin{teorem}[Jensen Eşitsizliği]
\label{thm:jensen}
$\varphi : \mathbb{R} \to \mathbb{R}$ konveks fonksiyon ve $E[|X|] < \infty$ ise
\[
  \varphi(E[X]) \leq E[\varphi(X)].
\]
\end{teorem}

\begin{proof}
$\varphi$ konveks $\Rightarrow$ her $\mu = E[X]$ için bir $c \in \mathbb{R}$ var:
$\varphi(x) \geq \varphi(\mu) + c(x - \mu)$ (destek doğrusu).
Beklenti alınca: $E[\varphi(X)] \geq \varphi(\mu) + c(E[X] - \mu) = \varphi(\mu) = \varphi(E[X])$.
\end{proof}

\begin{karsiornek}[Jensen Sıkılığı]
$\varphi$ katı konveks ve $X$ sabit değil ise eşitsizlik kesin katıdır:
$\varphi(E[X]) < E[\varphi(X)]$.
Örnek: $\varphi(x) = x^2$, $X$ sabit olmayan $\Rightarrow$ $(E[X])^2 < E[X^2]$,
yani $\mathrm{Var}(X) = E[X^2] - (E[X])^2 > 0$.
\end{karsiornek}

\begin{ornek}[Jensen Uygulamaları]
\begin{itemize}
  \item $\varphi(x) = e^x$ konveks: $e^{E[X]} \leq E[e^X]$.
  \item $\varphi(x) = -\ln x$ konveks ($x>0$): $-\ln E[X] \leq E[-\ln X]$,
        yani $\ln E[X] \geq E[\ln X]$ (AM-GM eşitsizliğinin olasılıksal formu).
  \item $\varphi(x) = |x|^p$ ($p \geq 1$) konveks: $|E[X]|^p \leq E[|X|^p]$ (güç eşitsizliği).
\end{itemize}
\end{ornek}

\begin{mikroozet}
Doğrusallık + Monotonluk + Jensen: beklentinin üç temel aracı.
Markov: $P(X \geq a) \leq E[X]/a$. Chebyshev: $P(|X-\mu| \geq k\sigma) \leq 1/k^2$.
\end{mikroozet}

% =====================================================================
\section{Varyans ve Standart Sapma}
\label{sec:varyans}
% =====================================================================

\begin{tanim}[Varyans ve Standart Sapma]
\label{def:varyans}
$E[X^2] < \infty$ ise $X$'in \textbf{varyansı} (variance):
\[
  \mathrm{Var}(X) = E[(X - E[X])^2] = E[X^2] - (E[X])^2.
\]
\textbf{Standart sapma:} $\sigma_X = \sqrt{\mathrm{Var}(X)}$.
\end{tanim}

\begin{proof}[Hesaplama Formülünün İspat]
$E[(X - \mu)^2] = E[X^2 - 2\mu X + \mu^2]
= E[X^2] - 2\mu E[X] + \mu^2 = E[X^2] - \mu^2$.
\end{proof}

\begin{onerme}[Varyansın Özellikleri]
\label{prop:varyans_ozellik}
\begin{enumerate}[(i)]
  \item $\mathrm{Var}(X) \geq 0$; $\mathrm{Var}(X) = 0 \Leftrightarrow X = \mu$ n.k.
  \item $\mathrm{Var}(aX + b) = a^2 \mathrm{Var}(X)$.
  \item $X, Y$ bağımsız ise $\mathrm{Var}(X + Y) = \mathrm{Var}(X) + \mathrm{Var}(Y)$.
  \item Genel: $\mathrm{Var}(X + Y) = \mathrm{Var}(X) + \mathrm{Var}(Y) + 2\mathrm{Cov}(X, Y)$.
\end{enumerate}
\end{onerme}

\begin{proof}
\textit{(ii)}: $E[(aX+b-aE[X]-b)^2] = a^2 E[(X-E[X])^2]$.\\
\textit{(iii)}: Bağımsız $\Rightarrow$ $\mathrm{Cov}(X,Y) = 0$ (aşağıda ispat edilecek).\\
\textit{(iv)}: $\mathrm{Var}(X+Y) = E[(X+Y-\mu_X-\mu_Y)^2]$; açılımda $E[(X-\mu_X)(Y-\mu_Y)] = \mathrm{Cov}(X,Y)$.
\end{proof}

\begin{ornek}[Binom Dağılımı]
$X \sim \Binom(n, p)$: $X = \sum_{i=1}^n X_i$ ($X_i \sim \mathrm{Bernoulli}(p)$ bağımsız).
$E[X_i] = p$, $\mathrm{Var}(X_i) = p(1-p)$.
Doğrusallık ve bağımsızlıktan: $E[X] = np$, $\mathrm{Var}(X) = np(1-p)$.
\end{ornek}

% =====================================================================
\section{Kovaryans ve Korelasyon}
\label{sec:kovaryans}
% =====================================================================

\begin{tanim}[Kovaryans ve Korelasyon]
\label{def:kovaryans}
$E[X^2], E[Y^2] < \infty$ ise
\begin{align*}
  \mathrm{Cov}(X, Y) &= E[(X - E[X])(Y - E[Y])] = E[XY] - E[X]E[Y], \\
  \mathrm{Cor}(X, Y) &= \rho_{XY} = \frac{\mathrm{Cov}(X,Y)}{\sqrt{\mathrm{Var}(X)\,\mathrm{Var}(Y)}}.
\end{align*}
\end{tanim}

\begin{teorem}[Korelasyonun Sınırları]
\label{thm:cauchy_schwarz}
$-1 \leq \mathrm{Cor}(X, Y) \leq 1$.
$|\rho_{XY}| = 1$ ancak ve ancak $Y = aX + b$ (n.k.) ise ($a \neq 0$).
\end{teorem}

\begin{proof}
\textbf{Cauchy-Schwarz eşitsizliği:} $(\mathrm{Cov}(X,Y))^2 \leq \mathrm{Var}(X) \cdot \mathrm{Var}(Y)$.
İspat: $f(t) = \mathrm{Var}(tX + Y) = t^2 \mathrm{Var}(X) + 2t\,\mathrm{Cov}(X,Y) + \mathrm{Var}(Y) \geq 0$
(diskriminant $\leq 0$): $4(\mathrm{Cov}(X,Y))^2 - 4\mathrm{Var}(X)\mathrm{Var}(Y) \leq 0$.
Eşitlik $f(t^*) = 0$ için; bu $t^*X + Y = 0$ n.k.\ anlamına gelir, yani $Y = -t^*X$ lineer ilişkisi.
\end{proof}

\begin{hataanalizi}
"$\rho = 0$ ise $X, Y$ bağımsızdır" — yanlış.
Bölüm~\ref{sec:rd_bagimsizlik}'daki $X \sim U(-1,1)$, $Y = X^2$ örneği:
$\mathrm{Cov}(X,Y) = E[X^3] - E[X]E[X^2] = 0 - 0 = 0$ ama bağımsız değiller.
Korelasyon yalnızca doğrusal ilişkiyi ölçer.
\end{hataanalizi}

% =====================================================================
\section{Momentler ve Merkezi Momentler}
\label{sec:momentler}
% =====================================================================

\begin{tanim}[Momentler]
\label{def:moment}
$E[|X|^r] < \infty$ ise ($r > 0$):
\begin{itemize}
  \item \textbf{$r$'inci moment:} $\mu_r' = E[X^r]$.
  \item \textbf{$r$'inci merkezi moment:} $\mu_r = E[(X-\mu)^r]$ ($\mu = E[X]$).
  \item \textbf{Standardize $r$'inci moment:} $\tilde{\mu}_r = \mu_r / \sigma^r$.
\end{itemize}
Önemli özel durumlar: $\mu_1' = E[X]$ (birinci moment), $\mu_2 = \mathrm{Var}(X)$.
\end{tanim}

\begin{tanim}[Çarpıklık ve Basıklık]
\label{def:carpiklik}
$\sigma > 0$ ve $E[|X|^4] < \infty$ olmak üzere:
\begin{align*}
  \text{Çarpıklık (skewness):} &\quad \gamma_1 = \tilde{\mu}_3 = \frac{E[(X-\mu)^3]}{\sigma^3}. \\
  \text{Basıklık (kurtosis):} &\quad \gamma_2 = \tilde{\mu}_4 - 3 = \frac{E[(X-\mu)^4]}{\sigma^4} - 3.
\end{align*}
$\gamma_2 = 0$ normal dağılımın referans değeridir (ekses basıklık).
$\gamma_2 > 0$: ağır kuyruk (leptokurtik); $\gamma_2 < 0$: hafif kuyruk (platikurtik).
\end{tanim}

\begin{ornek}[Normal Dağılımın Momentleri]
$X \sim \Normal(\mu, \sigma^2)$:
Simetri $\Rightarrow$ tüm tek merkezi momentler sıfır ($\gamma_1 = 0$).
$E[(X-\mu)^4] = 3\sigma^4$ $\Rightarrow$ $\gamma_2 = 0$.
Genel: $E[(X-\mu)^{2k}] = (2k-1)!! \cdot \sigma^{2k}$ (çift momentler formülü).
\end{ornek}

% =====================================================================
\section{Moment Üreten Fonksiyon}
\label{sec:muf}
% =====================================================================

\begin{kesif}
Bir rasgele değişkenin tüm momentlerini ($E[X], E[X^2], E[X^3], \ldots$) tek bir
fonksiyona kodlayabilir miyiz? Ve bu fonksiyon dağılımı tam olarak belirler mi?
\end{kesif}

\begin{tanim}[Moment Üreten Fonksiyon]
\label{def:muf}
$X$ rasgele değişkeninin \textbf{moment üreten fonksiyonu} (MÜF, moment generating function):
\[
  M_X(t) = E[e^{tX}], \qquad t \in \mathbb{R},
\]
eğer bu beklenti $t$'nin bir $(-h, h)$ ($h > 0$) komşuluğunda sonlu ise.
Bu koşul sağlanmıyorsa MÜF tanımsız sayılır.
\end{tanim}

\begin{teorem}[MÜF'ten Momentler]
\label{thm:muf_moment}
$M_X(t)$ açık bir $(- h, h)$ aralığında sonlu ise:
\begin{enumerate}[(i)]
  \item $E[|X|^r] < \infty$ her $r \geq 1$ için.
  \item $M_X(t) = \sum_{r=0}^\infty \frac{E[X^r]}{r!} t^r$ ($|t| < h$).
  \item $E[X^r] = M_X^{(r)}(0) = \left.\frac{d^r}{dt^r} M_X(t)\right|_{t=0}$.
\end{enumerate}
\end{teorem}

\begin{proof}
\textit{(ii)--(iii):} $e^{tX} = \sum_{r=0}^\infty \frac{(tX)^r}{r!}$.
$|t| < h$ için DCT (Bölüm~\ref{chp:olcum_teorisi}) uygulanabilir;
seri ve integral sıraları değiştirilebilir:
\[
  M_X(t) = E\!\left[\sum_{r=0}^\infty \frac{(tX)^r}{r!}\right]
  = \sum_{r=0}^\infty \frac{E[X^r]}{r!} t^r.
\]
$t = 0$'da $r$'inci türev: $M_X^{(r)}(0) = E[X^r]$.
\end{proof}

\begin{teorem}[MÜF Dağılımı Belirler]
\label{thm:muf_tek}
$M_X(t) = M_Y(t)$ açık bir aralıkta ise $X \stackrel{d}{=} Y$ (dağılımca eşit).
\end{teorem}

Bu teoremin ispatı için moment probleminin çözülebilirliği (Hamburger, Stieltjes)
ve Laplace dönüşümünün tek-türlülüğü gerekmektedir; ayrıntı için Billingsley~\cite{Billingsley1995}
veya Klenke~\cite{Klenke2014}'e bakılabilir.

\begin{teorem}[MÜF ve Bağımsızlık]
\label{thm:muf_bagimsiz}
$X, Y$ bağımsız ise $M_{X+Y}(t) = M_X(t) \cdot M_Y(t)$.
\end{teorem}

\begin{proof}
$M_{X+Y}(t) = E[e^{t(X+Y)}] = E[e^{tX} e^{tY}] = E[e^{tX}] \cdot E[e^{tY}] = M_X(t) M_Y(t)$.
(Bağımsızlık: $E[fg] = E[f]E[g]$.)
\end{proof}

\begin{ornek}[Normal Dağılımın MÜF'ü]
$X \sim \Normal(\mu, \sigma^2)$:
\[
  M_X(t) = \exp\!\left(\mu t + \tfrac{1}{2}\sigma^2 t^2\right), \qquad t \in \mathbb{R}.
\]
\textit{İspat:} $M_X(t) = \int_{-\infty}^{+\infty} e^{tx} \frac{1}{\sigma\sqrt{2\pi}} e^{-(x-\mu)^2/(2\sigma^2)} dx$.
Üstel ifadeler birleştirilir:
$tx - \frac{(x-\mu)^2}{2\sigma^2} = -\frac{(x-\mu-\sigma^2 t)^2}{2\sigma^2} + \mu t + \frac{\sigma^2 t^2}{2}$.
Gauss integrali 1'e tamamlanır; $M_X(t) = e^{\mu t + \sigma^2 t^2/2}$.\\
Bağımsız $X \sim \Normal(\mu_1, \sigma_1^2)$, $Y \sim \Normal(\mu_2, \sigma_2^2)$:
$M_{X+Y}(t) = e^{(\mu_1+\mu_2)t + (\sigma_1^2+\sigma_2^2)t^2/2}$,
yani $X + Y \sim \Normal(\mu_1+\mu_2, \sigma_1^2+\sigma_2^2)$.
\end{ornek}

\begin{hataanalizi}
"Her rasgele değişkenin MÜF'ü vardır" — yanlış.
Cauchy dağılımı için $E[e^{tX}] = \infty$ her $t \neq 0$ için.
Lognormal dağılımı için tüm momentler mevcut ama MÜF sonsuz ($t > 0$).
MÜF yokluğu problemi, karakteristik fonksiyon kullanılarak aşılır (bir sonraki kesim).
\end{hataanalizi}

\begin{mikroozet}
$M_X(t) = E[e^{tX}]$; $E[X^r] = M_X^{(r)}(0)$.
Bağımsız $\Rightarrow$ $M_{X+Y} = M_X \cdot M_Y$.
MÜF her dağılım için var olmayabilir.
\end{mikroozet}

% =====================================================================
\section{Karakteristik Fonksiyon}
\label{sec:kf}
% =====================================================================

\begin{tanimgerekce}
MÜF $E[e^{tX}]$ ($t$ gerçel) her zaman mevcut değildir.
$e^{itX}$ ($i = \sqrt{-1}$, $t$ gerçel) ise $|e^{itX}| = 1$ olduğundan
her zaman integrallenebilir: $E[|e^{itX}|] = 1 < \infty$.
Kompleks üsseli beklenti her dağılım için tanımlıdır.
\end{tanimgerekce}

\begin{tanim}[Karakteristik Fonksiyon]
\label{def:kf}
$X$ rasgele değişkeninin \textbf{karakteristik fonksiyonu} (KF, characteristic function):
\[
  \varphi_X(t) = E[e^{itX}] = E[\cos(tX)] + i \, E[\sin(tX)], \qquad t \in \mathbb{R}.
\]
\end{tanim}

\begin{teorem}[Karakteristik Fonksiyonun Temel Özellikleri]
\label{thm:kf_ozellik}
\begin{enumerate}[(i)]
  \item $\varphi_X(0) = 1$, $|\varphi_X(t)| \leq 1$.
  \item $\varphi_X$ düzgün süreklidir (uniformly continuous).
  \item $\varphi_X$ dağılımı tek türlü belirler (ters dönüşüm teoremi).
  \item $E[|X|^r] < \infty$ ise $\varphi_X^{(r)}(0) = i^r E[X^r]$.
  \item $X, Y$ bağımsız ise $\varphi_{X+Y}(t) = \varphi_X(t) \cdot \varphi_Y(t)$.
\end{enumerate}
\end{teorem}

\begin{proof}
\textit{(i)}: $|e^{itX}| = 1$, dolayısıyla $|\varphi_X(t)| = |E[e^{itX}]| \leq E[|e^{itX}|] = 1$.\\
\textit{(ii)}: $|\varphi_X(t+h) - \varphi_X(t)| = |E[e^{itX}(e^{ihX}-1)]| \leq E[|e^{ihX}-1|]$.
$h \to 0$: $|e^{ihX}-1| \leq 2$, DCT ile $E[|e^{ihX}-1|] \to 0$.\\
\textit{(v)}: MÜF için aynı argüman; $|e^{itX}|=1$ dolayısıyla sonluluk kaygısı yok.
\end{proof}

\begin{teorem}[Ters Dönüşüm]
\label{thm:kf_ters}
$X$ mutlak sürekli, $\int |\varphi_X(t)| \, dt < \infty$ ise
\[
  f_X(x) = \frac{1}{2\pi} \int_{-\infty}^{+\infty} e^{-itx} \varphi_X(t) \, dt.
\]
\end{teorem}

Bu, Fourier dönüşümünün ters formülüdür; gerçel analize (Plancherel teoremi) dayanır.

\begin{ornek}[Normal Dağılımın KF'si]
$X \sim \Normal(\mu, \sigma^2)$:
\[
  \varphi_X(t) = \exp\!\left(i\mu t - \tfrac{1}{2}\sigma^2 t^2\right).
\]
MÜF'te $t \to it$ ikamesi yapar gibi: $M_X(it) = e^{i\mu t - \sigma^2 t^2/2}$.
(Gerçek ispat için analitik sürdürüm veya doğrudan Gauss hesabı gereklidir.)
\end{ornek}

\begin{onerme}[Dağılımla Yakınsama ve KF]
\label{prop:kf_yaklasim}
$X_n \xrightarrow{d} X$ ancak ve ancak her $t \in \mathbb{R}$ için
$\varphi_{X_n}(t) \to \varphi_X(t)$ ise (Lévy süreklilik teoremi).
\end{onerme}

Bu önerme, Bölüm~\ref{chp:merkezi_limit}'teki Merkezi Limit Teoremi'nin kanıtında
kilit rol oynayacaktır.

\begin{mikroozet}
$\varphi_X(t) = E[e^{itX}]$: her dağılım için mevcut, düzgün sürekli, $|\varphi| \leq 1$.
Bağımsız toplam: $\varphi_{X+Y} = \varphi_X \cdot \varphi_Y$.
Dağılımla yakınsama $\Leftrightarrow$ KF'lerin noktasal yakınsaması (Lévy).
\end{mikroozet}

% =====================================================================
\section{Kümülan Üreten Fonksiyon}
\label{sec:kgf}
% =====================================================================

\begin{tanim}[Kümülan Üreten Fonksiyon]
\label{def:kgf}
$M_X(t)$ mevcut ise \textbf{kümülan üreten fonksiyon} (cumulant generating function):
\[
  K_X(t) = \ln M_X(t) = \sum_{r=1}^\infty \frac{\kappa_r}{r!} t^r.
\]
$\kappa_r = K_X^{(r)}(0)$: $r$'inci \textbf{kümülan} (cumulant).
\end{tanim}

\begin{onerme}[Önemli Kümanlar]
\label{prop:kumulanlar}
\begin{align*}
  \kappa_1 &= E[X] = \mu, \\
  \kappa_2 &= \mathrm{Var}(X) = \sigma^2, \\
  \kappa_3 &= E[(X-\mu)^3] \quad \text{(çarpıklıkla orantılı)}, \\
  \kappa_4 &= E[(X-\mu)^4] - 3\sigma^4 \quad \text{(ekses basıklıkla orantılı)}.
\end{align*}
Bağımsız $X, Y$ için $K_{X+Y}(t) = K_X(t) + K_Y(t)$, dolayısıyla
kümanlar toplanır: $\kappa_r(X+Y) = \kappa_r(X) + \kappa_r(Y)$.
\end{onerme}

\begin{proof}
$K_X(t) = \ln M_X(t)$; $K_X'(0) = M_X'(0)/M_X(0) = E[X]/1 = \mu$.
$K_X''(0) = E[X^2] - (E[X])^2 = \mathrm{Var}(X)$.
Bağımsız toplam: $\ln M_{X+Y} = \ln M_X + \ln M_Y$.
\end{proof}

\begin{ornek}[Normal Dağılımın Kümanları]
$X \sim \Normal(\mu, \sigma^2)$: $K_X(t) = \mu t + \sigma^2 t^2/2$.
Sadece $\kappa_1 = \mu$ ve $\kappa_2 = \sigma^2$ sıfırdan farklı;
$r \geq 3$ için $\kappa_r = 0$.
Bu, normal dağılımın kümülan diliyle karakterizasyonudur.
\end{ornek}

% =====================================================================
\section{Koşullu Beklenti}
\label{sec:kosullu_beklenti}
% =====================================================================

\begin{kesif}
$Y = 2X + \varepsilon$ ($\varepsilon \perp X$) modelini düşünün.
$X$'in değerini bildiğimizde $Y$ için en iyi tahminimiz nedir?
Sezgisel olarak $E[Y \mid X = x] = 2x$ diyoruz.
Peki ya $X$ yerine bir olay $\mathcal{G}$ verildiğinde?
"En iyi tahmin" kavramını nasıl soyutlarız?
\end{kesif}

\begin{tanim}[Koşullu Beklenti — Elemanter]
\label{def:kosullu_beklenti_elem}
\textbf{Kesikli:} $P(Y = y_j) > 0$ olayları üretilen $\mathcal{G} = \sigma(Y)$ için
\[
  E[X \mid Y = y_j] = \frac{E[X \cdot \mathbf{1}_{\{Y=y_j\}}]}{P(Y = y_j)}.
\]
\textbf{Sürekli:} $f_Y(y) > 0$ ise $E[X \mid Y = y] = \int x \, f_{X|Y}(x \mid y) \, dx$.
\end{tanim}

\begin{tanimgerekce}
Genel $\sigma$-cebiri için yukarıdaki formüller yetersiz kalır (sıfır olasılıklı koşullamada anlam kaybolur).
Radon-Nikodym teoremi, koşullu beklentiyi "en iyi $\mathcal{G}$-ölçülebilir tahmin"
olarak $L^2$-projeksiyon ile tanımlamamızı sağlar.
\end{tanimgerekce}

\begin{tanim}[Koşullu Beklenti — Genel]
\label{def:kosullu_beklenti}
$(\Omega, \mathcal{F}, P)$ üzerinde $X \in L^1$ ve $\mathcal{G} \subseteq \mathcal{F}$
alt $\sigma$-cebiri verilsin. $E[X \mid \mathcal{G}]$, şu iki koşulu sağlayan
(n.k.\ tek türlü) rasgele değişkendir:
\begin{enumerate}[(i)]
  \item $E[X \mid \mathcal{G}]$ \textbf{$\mathcal{G}$-ölçülebilir}.
  \item Her $G \in \mathcal{G}$ için $\int_G E[X \mid \mathcal{G}] \, dP = \int_G X \, dP$.
\end{enumerate}
Özel olarak $\mathcal{G} = \sigma(Y)$ için $E[X \mid Y] := E[X \mid \sigma(Y)]$.
\end{tanim}

\begin{proof}[Varlık ve Tek Türlülük]
Tek türlülük: $Z_1, Z_2$ her ikisi de koşulları sağlasın; $G = \{Z_1 > Z_2\} \in \mathcal{G}$:
$\int_G (Z_1 - Z_2) \, dP = 0$ ve $Z_1 - Z_2 > 0$ burada; yani $P(G) = 0$.
Benzer şekilde $\{Z_1 < Z_2\}$ için; dolayısıyla $Z_1 = Z_2$ n.k.\\
Varlık: $\nu(G) = \int_G X \, dP$ ($G \in \mathcal{G}$) $\sigma$-sonlu imzalı ölçü;
$P|_{\mathcal{G}}$'ye mutlak sürekli. Radon-Nikodym: $\nu = Z \cdot P|_{\mathcal{G}}$
yazan $\mathcal{G}$-ölçülebilir $Z$ var; $Z = E[X \mid \mathcal{G}]$.
\end{proof}

\begin{teorem}[Koşullu Beklentinin Temel Özellikleri]
\label{thm:kosullu_beklenti}
$X, Y \in L^1$, $\mathcal{G} \subseteq \mathcal{H} \subseteq \mathcal{F}$ alt $\sigma$-cebirleri.
\begin{enumerate}[(i)]
  \item \textbf{Doğrusallık:} $E[aX + bY \mid \mathcal{G}] = aE[X\mid\mathcal{G}] + bE[Y\mid\mathcal{G}]$.
  \item \textbf{Kule (yineleme) özelliği:} $E[E[X \mid \mathcal{H}] \mid \mathcal{G}] = E[X \mid \mathcal{G}]$.
        Özellikle $E[E[X \mid Y]] = E[X]$ (toplam beklenti formülü).
  \item \textbf{$\mathcal{G}$-ölçülebilir çarpan:} $Z$ $\mathcal{G}$-ölçülebilir ise
        $E[ZX \mid \mathcal{G}] = Z \cdot E[X \mid \mathcal{G}]$.
  \item \textbf{Bağımsızlık:} $X$ $\mathcal{G}$'den bağımsız ise $E[X \mid \mathcal{G}] = E[X]$.
  \item \textbf{Monotonluk:} $X \leq Y$ n.k.\ ise $E[X \mid \mathcal{G}] \leq E[Y \mid \mathcal{G}]$ n.k.
  \item \textbf{Jensen:} $\varphi$ konveks ise $\varphi(E[X \mid \mathcal{G}]) \leq E[\varphi(X) \mid \mathcal{G}]$.
\end{enumerate}
\end{teorem}

\begin{proof}
\textit{(ii) Kule özelliği:} $\mathcal{G} \subseteq \mathcal{H}$.
Her $G \in \mathcal{G} \subseteq \mathcal{H}$ için tanım koşulu (ii):
$\int_G E[X\mid\mathcal{H}] \, dP = \int_G X \, dP$.
Ama $E[X\mid\mathcal{H}]$'nın $\mathcal{G}$-koşullu beklentisi için de bu değer kullanılabilir;
dolayısıyla $E[E[X\mid\mathcal{H}]\mid\mathcal{G}]$ tanım koşullarını $E[X\mid\mathcal{G}]$ ile aynı şekilde sağlar.
Tek türlülükten eşitlik.\\
\textit{(iii):} $G \in \mathcal{G}$: $\int_G ZE[X\mid\mathcal{G}] \, dP = \int_G Z \cdot E[X\mid\mathcal{G}] \, dP$.
$Z$ basamak fonksiyonu ise doğrudan; genel $Z$ için yaklaşım.
\end{proof}

\begin{onerme}[$L^2$ Projeksiyon Yorumu]
\label{prop:l2_projeksiyon}
$X \in L^2(\Omega, \mathcal{F}, P)$ için $E[X \mid \mathcal{G}]$,
$L^2(\Omega, \mathcal{G}, P)$ alt uzayına $X$'in dik projeksiyonudur:
\[
  E[X \mid \mathcal{G}] = \arg\min_{Z \in L^2(\mathcal{G})} E[(X - Z)^2].
\]
\end{onerme}

\begin{proof}
$Z \in L^2(\mathcal{G})$ için:
$E[(X-Z)^2] = E[(X - E[X\mid\mathcal{G}] + E[X\mid\mathcal{G}] - Z)^2]$.
Açılımda çapraz terim $E[(X - E[X\mid\mathcal{G}])(E[X\mid\mathcal{G}] - Z)]$:
$(E[X\mid\mathcal{G}] - Z)$ $\mathcal{G}$-ölçülebilir; (iii) özelliğiyle çapraz terim sıfır.
Dolayısıyla $E[(X-Z)^2] = E[(X-E[X\mid\mathcal{G}])^2] + E[(E[X\mid\mathcal{G}]-Z)^2]$,
sağ tarafın minimumu $Z = E[X \mid \mathcal{G}]$'de sağlanır.
\end{proof}

\begin{ornek}[Regresyon]
$E[Y \mid X] = g(X)$; burada $g(x) = E[Y \mid X = x]$.
$E[Y \mid X]$ koşullu beklentisi, $Y$'yi $X$ bilgisiyle tahmin etmede
en küçük ortalama karesel hata veren tahminci (MMSE tahmin edici) olup
"doğrusal olmayan regresyon" kavramıyla özdeşleşir.
Doğrusal kısıt altında ($g(X) = a + bX$): $b = \mathrm{Cov}(X,Y)/\mathrm{Var}(X)$.
\end{ornek}

\begin{kendinleifade}
Koşullu beklentinin "bilgiyi güncellemek" yorumunu kendi cümlelerinizle açıklayın.
"Kule özelliği" ne anlama gelir? Somut bir örnek üretin.
\ogrencialani[2.5cm]
\end{kendinleifade}

% ---- Disiplinlerarası Bağlantı (TYMM) ----
\begin{kopru}[Disiplinlerarası — Ekonomi, Fizik, Bilgisayar Bilimi]
\textbf{Ekonomi:} Beklenen fayda teorisi (von Neumann-Morgenstern), karar teorisinin
matematiksel temelidir; Jensen eşitsizliği risk kaçınımını açıklar ($u$ içbükey $\Rightarrow$ $u(E[X]) > E[u(X)]$).\\
\textbf{Makine öğrenmesi:} Koşullu beklenti $E[Y|X]$, gözetimli öğrenmede kayıp fonksiyonu
minimize eden tahmin edicinin kendisidir; nöral ağlar bu fonksiyonu yaklaşık öğrenir.\\
\textbf{Kuantum mekaniği:} Observable $A$ operatörünün beklenti değeri $\langle A \rangle = \langle \psi | A | \psi \rangle$,
istatistiksel beklentiyle yapısal olarak örtüşür.
\end{kopru}

% ---- Öz-Yansıtma ----
\begin{ozyansima}
\begin{itemize}[leftmargin=1.8em]
  \item Beklentinin neden Lebesgue integrali olarak tanımlanmasının gerekli olduğunu açıklayabilir misiniz?
  \item Jensen eşitsizliği sizce en güçlü hangi uygulamada karşınıza çıktı?
  \item MÜF ile KF arasındaki temel fark nedir, ve hangisi daha "evrensel"dir?
  \item Koşullu beklentiyi $L^2$ projeksiyonu olarak görmek, sezginizi nasıl değiştirdi?
\end{itemize}
\end{ozyansima}

% ---- Köprü Notu ----
\begin{kopru}
\textbf{Önceki bölüm:} Bölüm~\ref{chp:rasgele_degiskenler} — dağılım fonksiyonu ve PDF;
burada bunlar Lebesgue integrali ile beklentiye dönüştü.\\
\textbf{Sonraki bölüm:} Bölüm~\ref{chp:dagilim_aileleri} — somut dağılımların
beklenti, varyans, MÜF hesapları bu bölümün araçlarıyla yapılacak.\\
\textbf{İleride:} Bölüm~\ref{chp:buyuk_sayilar}'da Chebyshev ile ZBSK,
Bölüm~\ref{chp:merkezi_limit}'te KF ile MLT kanıtı.
\defterbaglantisi{bolum04\_beklenti.ipynb}{%
  MÜF hesabı, Jensen görselleştirmesi, koşullu beklenti simülasyonu}
\end{kopru}

% =====================================================================
\section{Olasılık Eşitsizlikleri}
\label{sec:olasilik_esitsizlikleri}
% =====================================================================

\begin{motivasyon}[Neden Eşitsizlikler?]
Bir dağılımı tam olarak bilmeden $P(|X - \mu| > \varepsilon)$ hakkında ne söyleyebiliriz?
Bu soru yalnızca momentlere dayalı sınırlarla yanıtlanabilir.
Markov, Chebyshev ve Jensen eşitsizlikleri bu sınırların sistematik kaynağıdır;
Büyük Sayılar Kanunu ve Büyük Sapmalar teorisinin altyapısını oluştururlar.
\end{motivasyon}

\begin{teorem}[Markov Eşitsizliği]
\label{thm:markov_beklenti}
$X \geq 0$ ve $E[X] < \infty$. Her $a > 0$ için
\[
  P(X \geq a) \leq \frac{E[X]}{a}.
\]
\end{teorem}

\begin{proof}
$X \geq a\,\mathbf{1}_{\{X \geq a\}}$ (çünkü $X \geq 0$).
Beklenti alınca: $E[X] \geq a\,P(X \geq a)$. Düzenlenince sonuç.
\end{proof}

\begin{teorem}[Chebyshev Eşitsizliği — Genelleştirilmiş]
\label{thm:chebyshev_genellestirilmis}
$\varphi : \mathbb{R} \to [0, \infty)$ simetrik ve artan ($|x|$'de) fonksiyon,
$\varphi(a) > 0$. O zaman
\[
  P(|X - \mu| \geq a) \leq \frac{E[\varphi(X-\mu)]}{\varphi(a)}.
\]
Özel durum $\varphi(x) = x^2$: $P(|X-\mu|\geq a) \leq \mathrm{Var}(X)/a^2$.
\end{teorem}

\begin{proof}
$\varphi(|X-\mu|) \geq \varphi(a)\,\mathbf{1}_{\{|X-\mu| \geq a\}}$ ($\varphi$ artan).
Markov: $E[\varphi(|X-\mu|)] \geq \varphi(a)\,P(|X-\mu|\geq a)$.
\end{proof}

\begin{teorem}[Tek Taraflı Chebyshev (Cantelli Eşitsizliği)]
\label{thm:cantelli}
$E[X] = \mu$, $\mathrm{Var}(X) = \sigma^2 < \infty$. Her $\lambda > 0$ için
\[
  P(X - \mu \geq \lambda) \leq \frac{\sigma^2}{\sigma^2 + \lambda^2}.
\]
\end{teorem}

\begin{proof}
Her $t > 0$ için $P(X-\mu \geq \lambda) = P(X-\mu+t \geq \lambda+t)$.
Markov ($X-\mu+t \geq 0$ değil; bunun yerine $(X-\mu+t)^2$ ile):
$P((X-\mu+t)^2 \geq (\lambda+t)^2) \leq E[(X-\mu+t)^2]/(\lambda+t)^2 = (\sigma^2+t^2)/(\lambda+t)^2$.
$t = \sigma^2/\lambda$'da minimize et: $(\sigma^2 + \sigma^4/\lambda^2)/(\lambda + \sigma^2/\lambda)^2 = \sigma^2/(\sigma^2+\lambda^2)$.
\end{proof}

\begin{ornek}[Eşitsizliklerin Karşılaştırması]
$X \sim \Normal(0,1)$, $a = 3$.
\begin{itemize}
  \item Markov ($|X|$ ile): $P(|X| \geq 3) \leq E[|X|]/3 = \sqrt{2/\pi}/3 \approx 0.266$.
  \item Chebyshev: $P(|X| \geq 3) \leq 1/9 \approx 0.111$.
  \item Cantelli (tek taraf): $P(X \geq 3) \leq 1/(1+9) = 0.1$.
  \item Gerçek değer: $P(|X| \geq 3) = 2(1-\Phi(3)) \approx 0.0027$.
\end{itemize}
Eşitsizlikler "kötü" ama dağılım bilgisi olmadan geçerlidir.
\end{ornek}

\begin{teorem}[Hoeffding Eşitsizliği]
\label{thm:hoeffding}
$X_1, \ldots, X_n$ bağımsız, $a_i \leq X_i \leq b_i$ n.k., $E[X_i] = \mu_i$.
Her $t > 0$ için
\[
  P\!\left(\sum_{i=1}^n (X_i - \mu_i) \geq t\right) \leq \exp\!\left(-\frac{2t^2}{\sum_{i=1}^n(b_i-a_i)^2}\right).
\]
\end{teorem}

\begin{proof}[İspat Taslağı]
Hoeffding lemması: $E[X_i] = 0$, $a_i \leq X_i \leq b_i$ ise
$E[e^{sX_i}] \leq e^{s^2(b_i-a_i)^2/8}$.
(Konvekslik: $e^{sx} \leq \frac{b-x}{b-a}e^{sa} + \frac{x-a}{b-a}e^{sb}$;
beklenti al, $f(u) = -su + e^u$'nun Taylor açılımıyla sınırla.)
Markov + MÜF çarpımı: $P(\sum(X_i-\mu_i)\geq t) \leq e^{-st}\prod e^{s^2(b_i-a_i)^2/8}$.
$s = 4t/\sum(b_i-a_i)^2$'de minimize et: istenen sınır.
\end{proof}

\begin{ornek}[Makine Öğrenmesi — Genelleme Hatası]
$n$ bağımsız veri; her $X_i \in [0,1]$, $E[X_i] = \mu$.
Hoeffding: $P(\bar{X}_n - \mu \geq \varepsilon) \leq e^{-2n\varepsilon^2}$.
$n = 1000$, $\varepsilon = 0.05$: $P \leq e^{-5} \approx 0.007$.
PAC öğrenmesinde bu tür sınırlar "garantili öğrenme" teorisinin temelidir.
\end{ornek}

\begin{mikroozet}
Markov: $P(X \geq a) \leq E[X]/a$ ($X \geq 0$).
Chebyshev: $P(|X-\mu|\geq a) \leq \sigma^2/a^2$.
Cantelli: tek taraflı, $\sigma^2/(\sigma^2+\lambda^2)$.
Hoeffding: sınırlı değişkenler için üstel sınır.
\end{mikroozet}

% =====================================================================
\section{Değişim Katsayısı ve Standartlaştırma}
\label{sec:degisim_katsayisi}
% =====================================================================

\begin{tanim}[Değişim Katsayısı]
\label{def:degisim_katsayisi}
$\mu = E[X] \neq 0$ ve $\sigma = \sqrt{\mathrm{Var}(X)}$ için
\textbf{değişim katsayısı} (coefficient of variation):
\[
  \mathrm{CV}(X) = \frac{\sigma}{|\mu|}.
\]
Boyutsuz göreli değişkenlik ölçüsü; farklı birim veya ölçekteki verileri karşılaştırır.
\end{tanim}

\begin{ornek}[CV Karşılaştırması]
İki yatırım:
$X$: $\mu_X = 100$, $\sigma_X = 15$ $\Rightarrow$ $\mathrm{CV}(X) = 0.15$.
$Y$: $\mu_Y = 1000$, $\sigma_Y = 120$ $\Rightarrow$ $\mathrm{CV}(Y) = 0.12$.
Mutlak değişkenlik ($\sigma$) bakımından $Y$ daha değişken;
göreli değişkenlik (CV) bakımından $X$ daha riskli.
\end{ornek}

\begin{tanim}[Çarpıklık ve Basıklık]
\label{def:carpiklik_basiklik}
$\mu_k' = E[(X-\mu)^k]$ $k$'ıncı merkezi moment.
\[
  \text{Çarpıklık (skewness):} \quad \gamma_1 = \frac{\mu_3'}{\sigma^3}.
\]
\[
  \text{Basıklık (kurtosis):} \quad \gamma_2 = \frac{\mu_4'}{\sigma^4} - 3.
\]
$\gamma_2 = 0$: Normal dağılım. $\gamma_2 > 0$: ağır kuyruk (leptokurtik).
$\gamma_2 < 0$: hafif kuyruk (platikurtik).
\end{tanim}

\begin{teorem}[Normal Dağılımın Momentleri]
\label{thm:normal_momentler}
$X \sim \Normal(0,1)$: tüm tek momentler sıfır;
$E[X^{2k}] = (2k-1)!! = (2k)!/(2^k k!)$.
Özellikle: $E[X^4] = 3$ $\Rightarrow$ $\gamma_2 = 3 - 3 = 0$.
\end{teorem}

\begin{proof}
$E[X^{2k}] = \int_{-\infty}^\infty x^{2k} \phi(x) dx$.
Parçalı integral: $u = x^{2k-1}$, $dv = x\phi(x)dx = -\phi'(x)dx$; $v = -\phi(x)$.
$\int x^{2k}\phi(x)dx = [-x^{2k-1}\phi(x)]_{-\infty}^\infty + (2k-1)\int x^{2k-2}\phi(x)dx = (2k-1)E[X^{2k-2}]$.
Yinelemeli: $E[X^{2k}] = (2k-1)(2k-3)\cdots 1 = (2k-1)!!$.
\end{proof}

\begin{ornek}[Çarpıklık Yorumu]
$X \sim \Exp(\lambda)$: $\mu_3' = 2/\lambda^3$, $\sigma^3 = 1/\lambda^3$, $\gamma_1 = 2 > 0$: sağa çarpık.
$X \sim \Beta(2,5)$: $\gamma_1 = 2(5-2)\sqrt{(2+5+1)/((2\cdot5)\cdot 1)} / (2+5+2) > 0$: sağa çarpık.
Gelir dağılımları tipik olarak pozitif çarpık: kalabalık düşük-orta, seyrek yüksek gelirler.
\end{ornek}

% =====================================================================
\section{L$^p$ Uzayları ve Beklenti Eşitsizlikleri}
\label{sec:lp_uzaylari}
% =====================================================================

\begin{tanim}[$L^p$ normu]\label{def:lp_norm}
$p\geq1$. Rasgele değişken $X$ için
\[
  \|X\|_p = (\Beklenti[|X|^p])^{1/p}.
\]
$L^p(\Omega,\mathcal{F},P)=\{X:\|X\|_p<\infty\}$: $p$'inci mutlak moment sonlu olan rasgele değişkenler sınıfı.
\end{tanim}

\begin{teorem}[Lyapunov eşitsizliği]\label{thm:lyapunov_esit}
$1\leq r<s$: $\|X\|_r\leq\|X\|_s$.
Yani $L^s\subset L^r$ ve $r < s$ için $r$'inci moment $s$'inci momentten daha sık sonludur.
\end{teorem}

\begin{proof}
Jensen: $\phi(x)=x^{s/r}$ dışbükey ($s>r$). $\Beklenti[|X|^r]^{s/r}\leq\Beklenti[|X|^s]$. Her iki tarafın $r/s$'inci kuvveti: $\Beklenti[|X|^r]^{1}\leq\Beklenti[|X|^s]^{r/s}$, yani $\|X\|_r^r\leq\|X\|_s^r$, bu da $\|X\|_r\leq\|X\|_s$.
\end{proof}

\begin{teorem}[Hölder eşitsizliği]\label{thm:holder}
$1/p+1/q=1$, $p,q\geq1$:
\[
  \Beklenti[|XY|] \leq \|X\|_p\cdot\|Y\|_q.
\]
Özel durum $p=q=2$: Cauchy-Schwarz $\Beklenti[|XY|]\leq\sqrt{\Beklenti[X^2]\Beklenti[Y^2]}$.
\end{teorem}

\begin{proof}[Proof (Young ile)]
Young eşitsizliği: $ab\leq a^p/p+b^q/q$. $a=|X|/\|X\|_p$, $b=|Y|/\|Y\|_q$: $\frac{|XY|}{\|X\|_p\|Y\|_q}\leq\frac{|X|^p}{p\|X\|_p^p}+\frac{|Y|^q}{q\|Y\|_q^q}$. Beklenti alınca sağ taraf $1/p+1/q=1$.
\end{proof}

\begin{teorem}[Minkowski eşitsizliği]\label{thm:minkowski}
$p\geq1$: $\|X+Y\|_p\leq\|X\|_p+\|Y\|_p$ (üçgen eşitsizliği).
\end{teorem}

\begin{ornek}[$L^p$ dahil olma ilişkisi]\label{ex:lp_dahil}
Sonlu $\Omega=\{1,\ldots,n\}$ (her eleman eşit olasılıklı). $\|X\|_p=(\frac1n\sum|x_i|^p)^{1/p}$: RMS$(p=2)>$ortalama$(p=1)$ — karesel ortalamanın aritmetik ortalamayı aşması. $\|X\|_\infty=\max|x_i|$ = supremum norm.
\end{ornek}

% =====================================================================
\section{Koşullu Beklentinin Geometrisi}
\label{sec:kb_geometri}
% =====================================================================

\begin{tanim}[Projeksiyon yorumu]\label{def:projeksiyon}
$L^2(\Omega,\mathcal{F},P)$ Hilbert uzayı, iç çarpım $\langle X,Y\rangle=\Beklenti[XY]$. $\Beklenti[X|\mathcal{G}]$ (alt $\sigma$-cebir $\mathcal{G}$ üzerine koşullu beklenti), $X$'in $L^2(\Omega,\mathcal{G},P)$ altuzayına \emph{ortogonal projeksiyonudur}:
\[
  \Beklenti[\Beklenti[X|\mathcal{G}]Y] = \Beklenti[XY] \quad \forall Y\in L^2(\Omega,\mathcal{G},P).
\]
\end{tanim}

\begin{teorem}[Projeksiyon karakterizasyonu]\label{thm:projeksiyon}
$Z=\Beklenti[X|\mathcal{G}]$ iki koşulu sağlayan tek $\mathcal{G}$-ölçülebilir rasgele değişkendir:
\begin{enumerate}[(i)]
  \item $\Beklenti[ZY]=\Beklenti[XY]$ her $\mathcal{G}$-ölçülebilir $Y\in L^2$ için.
  \item $Z$, $\Beklenti[X-Z)^2]$'yi minimize eden — MMSE tahmincisi.
\end{enumerate}
\end{teorem}

\begin{proof}[Proof (taslak)]
(ii): $\Beklenti[(X-Z')^2]\geq\Beklenti[(X-Z)^2]$ her $\mathcal{G}$-ölçülebilir $Z'$ için. Genişlet: $\Beklenti[(X-Z')^2]=\Beklenti[(X-Z)^2]+\Beklenti[(Z-Z')^2]\geq\Beklenti[(X-Z)^2]$; orta terim sıfır (ortogonallik).
\end{proof}

\begin{ornek}[Lineer MMSE]\label{ex:lineer_mmse}
$(X,Y)$ ortak $L^2$. Lineer tahminci $\hat X = aY+b$ üzerinde MMSE:
$a^*=\Cov(X,Y)/\Var(Y)$, $b^*=E[X]-a^*E[Y]$.
$(X,Y)$ birlikte normal ise bu tahmin edici optimal (doğrusal olmayan gerekmez) ve $\hat X=\Beklenti[X|Y]$.
\end{ornek}

\begin{disiplinlerarasibaglan}
\textbf{Sinyal İşleme:} Wiener filtresi, gürültülü sinyalden MMSE tahmincisidir — koşullu beklentinin projeksiyon yorumunun mühendislik uygulaması.

\textbf{Ekonometri:} OLS (En Küçük Kareler) tahmincisi, $Y$'yi doğrusal $\mathbf{x}$ fonksiyonlarına projeksiyon olarak yorumlanır; $E[Y|\mathbf{X}=\mathbf{x}]$ doğrusal ise OLS MMSE elde eder.
\end{disiplinlerarasibaglan}

\begin{mikroozet}
CV: boyutsuz göreli değişkenlik; $\sigma/|\mu|$.
Çarpıklık: $\mu_3'/\sigma^3$; yönsel asimetri.
Basıklık fazlası: $\mu_4'/\sigma^4 - 3$; Normal'e göre kuyruk ağırlığı.
$L^p$ uzayları: Lyapunov ($\|X\|_r\leq\|X\|_s$, $r<s$), Hölder, Minkowski.
Koşullu beklenti = ortogonal projeksiyon $\Rightarrow$ MMSE tahmincisi.
\end{mikroozet}

% =====================================================================
%  ALIŞTIRMALAR
% =====================================================================

\cozumlualistirmalar

\begin{alistirma}[Al.4.1]{A}{Beklenti Hesabı}
$X$ PMF'si $p(k) = (1/2)^k$, $k = 1, 2, 3, \ldots$ olan kesikli rasgele değişken.
$E[X]$ ve $E[X^2]$'yi hesaplayın. Ardından $\mathrm{Var}(X)$'i bulun.
\end{alistirma}
\alcozum{%
$E[X] = \sum_{k=1}^\infty k (1/2)^k$. $\sum_{k=1}^\infty k r^k = r/(1-r)^2$ ($|r|<1$) ile $r=1/2$:
$E[X] = (1/2)/(1/2)^2 = 2$.\\
$E[X^2] = \sum_{k=1}^\infty k^2 (1/2)^k$. $\sum k^2 r^k = r(1+r)/(1-r)^3$: $E[X^2] = (1/2)(3/2)/(1/2)^3 = 6$.\\
$\mathrm{Var}(X) = 6 - 4 = 2$.}

\begin{alistirma}[Al.4.2]{A}{Jensen Eşitsizliği}
$X > 0$ rasgele değişken, $E[X] = 2$.
(a) $E[\ln X]$'i en fazla kaç yapabilirsiniz? Üst sınırı bulun.
(b) $E[1/X]$'in alt sınırını Jensen ile bulun.
\end{alistirma}
\alcozum{%
(a) $\varphi(x) = \ln x$ içbükey ($\varphi'' < 0$), yani $-\ln x$ konveks:
$-\ln E[X] \leq E[-\ln X]$, yani $E[\ln X] \leq \ln E[X] = \ln 2$.\\
(b) $\varphi(x) = 1/x$ konveks ($x>0$): Jensen $\Rightarrow$ $1/E[X] \leq E[1/X]$, yani $E[1/X] \geq 1/2$.}

\begin{alistirma}[Al.4.3]{B}{MÜF ile Moment Bulma}
$X \sim \mathrm{Pois}(\lambda)$: PMF $p(k) = e^{-\lambda}\lambda^k/k!$, $k=0,1,2,\ldots$
(a) $M_X(t) = e^{\lambda(e^t-1)}$ olduğunu gösterin.
(b) $M_X'(0)$ ve $M_X''(0)$ ile $E[X]$ ve $E[X^2]$'yi bulun.
(c) $\mathrm{Var}(X)$'i hesaplayın.
\end{alistirma}
\alcozum{%
(a) $M_X(t) = \sum_{k=0}^\infty e^{tk} e^{-\lambda} \lambda^k/k! = e^{-\lambda} \sum_{k=0}^\infty (\lambda e^t)^k/k!
= e^{-\lambda} e^{\lambda e^t} = e^{\lambda(e^t-1)}$.\\
(b) $M_X'(t) = \lambda e^t \cdot e^{\lambda(e^t-1)}$; $M_X'(0) = \lambda \cdot 1 = \lambda = E[X]$.
$M_X''(t) = (\lambda e^t + \lambda^2 e^{2t}) e^{\lambda(e^t-1)}$;
$M_X''(0) = \lambda + \lambda^2 = E[X^2]$.\\
(c) $\mathrm{Var}(X) = E[X^2] - (E[X])^2 = \lambda + \lambda^2 - \lambda^2 = \lambda$.
Poisson: beklenti = varyans = $\lambda$.}

\begin{alistirma}[Al.4.4]{B}{Koşullu Beklenti}
$X \sim \mathrm{Uniform}(0,1)$, $Y \mid X = x \sim \mathrm{Uniform}(0, x)$.
(a) $E[Y \mid X = x]$'i bulun.
(b) $E[Y]$'yi kule özelliğiyle hesaplayın.
(c) $\mathrm{Var}(Y)$'yi $\mathrm{Var}(Y) = E[\mathrm{Var}(Y\mid X)] + \mathrm{Var}(E[Y\mid X])$ ile bulun.
\end{alistirma}
\alcozum{%
(a) $E[Y \mid X=x] = x/2$.\\
(b) $E[Y] = E[E[Y\mid X]] = E[X/2] = (1/2) \cdot 1/2 = 1/4$.\\
(c) $\mathrm{Var}(Y\mid X=x) = x^2/12$.
$E[\mathrm{Var}(Y\mid X)] = E[X^2/12] = (1/3)/12 = 1/36$.
$\mathrm{Var}(E[Y\mid X]) = \mathrm{Var}(X/2) = (1/4)\mathrm{Var}(X) = (1/4)(1/12) = 1/48$.
$\mathrm{Var}(Y) = 1/36 + 1/48 = 4/144 + 3/144 = 7/144$.}

\begin{alistirma}[Al.4.4b]{B}{Karakteristik Fonksiyon ve Bağımsızlık}
$X,Y$ bağımsız rasgele değişkenler; $\varphi_X(t)=E[e^{itX}]$, $\varphi_Y(t)=E[e^{itY}]$.
(a) $\varphi_{X+Y}(t)=\varphi_X(t)\varphi_Y(t)$ olduğunu gösterin.
(b) $X,Y\iid\Normal(0,1)$: $\varphi_X(t)=e^{-t^2/2}$ kullanarak $Z=X+Y$ dağılımını bulun.
(c) $X\sim\Pois(\lambda_1)$, $Y\sim\Pois(\lambda_2)$ bağımsız: $Z=X+Y$ dağılımını KF yoluyla belirleyin.
\end{alistirma}
\alcozum{%
(a) $\varphi_{X+Y}(t)=E[e^{it(X+Y)}]=E[e^{itX}e^{itY}]=E[e^{itX}]E[e^{itY}]=\varphi_X(t)\varphi_Y(t)$.
(Bağımsızlık beklentinin çarpımını garanti eder.)

(b) $\varphi_Z(t)=\varphi_X(t)\varphi_Y(t)=e^{-t^2/2}\cdot e^{-t^2/2}=e^{-t^2}=e^{-(\sqrt{2})^2t^2/2}$.
Bu $\Normal(0,2)$'nin KF'sidir. Dolayısıyla $Z\sim\Normal(0,2)$.

(c) $\varphi_{\Pois(\lambda)}(t)=\exp(\lambda(e^{it}-1))$.
$\varphi_Z(t)=\exp(\lambda_1(e^{it}-1))\cdot\exp(\lambda_2(e^{it}-1))=\exp((\lambda_1+\lambda_2)(e^{it}-1))$.
Bu $\Pois(\lambda_1+\lambda_2)$'nin KF'sidir. $Z\sim\Pois(\lambda_1+\lambda_2)$. $\checkmark$}

\begin{alistirma}[Al.4.4c]{C}{Toplam Kuralı ile Ortak Dağılım Analizi}
$(X,Y)$ ortak yoğunluğu $f(x,y)=6xy^2$, $0<x<1$, $0<y<1$.
(a) Marjinal yoğunlukları $f_X(x)$ ve $f_Y(y)$'yi bulun.
(b) $X$ ve $Y$ bağımsız mı?
(c) $\Cov(X,Y)$, $\Cor(X,Y)$'yi hesaplayın.
(d) $E[Y\mid X=0.5]$'i bulun. Kule özelliğiyle $E[Y]$'yi doğrulayın.
\end{alistirma}
\alcozum{%
(a) $f_X(x)=\int_0^1 6xy^2\,dy=6x\cdot[y^3/3]_0^1=2x$, $x\in(0,1)$.\\
$f_Y(y)=\int_0^1 6xy^2\,dx=6y^2\cdot[x^2/2]_0^1=3y^2$, $y\in(0,1)$.

(b) $f(x,y)=6xy^2=(2x)(3y^2)=f_X(x)f_Y(y)$. Evet, $X$ ve $Y$ \textbf{bağımsız}.

(c) Bağımsız olduklarından $\Cov(X,Y)=0$ ve $\Cor(X,Y)=0$.
Doğrulama: $E[X]=\int_0^1 2x^2dx=2/3$; $E[X^2]=\int_0^1 2x^3dx=1/2$; $\Var(X)=1/2-4/9=1/18$.\\
$E[Y]=\int_0^1 3y^3dy=3/4$; $E[Y^2]=\int_0^1 3y^4dy=3/5$; $\Var(Y)=3/5-9/16=48/80-45/80=3/80$.\\
$E[XY]=E[X]E[Y]=(2/3)(3/4)=1/2$. $\Cov(X,Y)=E[XY]-E[X]E[Y]=1/2-1/2=0$. $\checkmark$

(d) $f_{Y|X}(y|x)=f(x,y)/f_X(x)=6xy^2/(2x)=3y^2$ — $X$'den bağımsız! $E[Y|X=0.5]=E[Y]=\int_0^1 3y^3dy=3/4$.\\
Kule: $E[E[Y|X]]=E[3/4]=3/4=E[Y]$. $\checkmark$}

% ---

\alistirmalar

\begin{alistirma}[Al.4.5]{A}{Beklenti Özellikleri}
$X$ rasgele değişken, $E[X] = 3$, $E[X^2] = 13$.
$E[(2X-1)^2]$ ve $\mathrm{Var}(3X+5)$'i hesaplayın.
\end{alistirma}
\alcozum{%
$E[(2X-1)^2] = 4E[X^2] - 4E[X] + 1 = 52 - 12 + 1 = 41$.
$\mathrm{Var}(X) = 13-9=4$. $\mathrm{Var}(3X+5) = 9\cdot 4 = 36$.}

\begin{alistirma}[Al.4.6]{A}{Chebyshev Uygulaması}
$X$ rasgele değişken, $E[X] = 50$, $\mathrm{Var}(X) = 25$.
(a) $P(|X - 50| \geq 10)$'un Chebyshev üst sınırını bulun.
(b) $P(40 \leq X \leq 60) \geq ?$ alt sınırını verin.
\end{alistirma}
\alcozum{%
(a) $\varepsilon = 10$: $P(|X-50|\geq 10) \leq 25/100 = 1/4$.
(b) $P(40 \leq X \leq 60) = 1 - P(|X-50| \geq 10) \geq 3/4$.}

\begin{alistirma}[Al.4.7]{B}{Markov Eşitsizliği ve Kesinlik}
$X \geq 0$, $E[X] = \mu$.
(a) $P(X \geq k\mu) \leq 1/k$ olduğunu gösterin.
(b) $P(X \geq 2\mu) = 1/2$ sağlayan bir dağılım inşa edin (Markov sıkılığı).
\end{alistirma}
\alcozum{%
(a) $a = k\mu$ ile Markov: $P(X \geq k\mu) \leq E[X]/(k\mu) = 1/k$.\\
(b) $P(X=0) = 1/2$, $P(X = 2\mu) = 1/2$: $E[X] = \mu$, $P(X \geq 2\mu) = 1/2$. Eşitlik sağlanır.}

\begin{alistirma}[Al.4.8]{B}{Kovaryans Matrisi}
$(X, Y, Z)$: $\mathrm{Var}(X) = 1$, $\mathrm{Var}(Y) = 4$, $\mathrm{Var}(Z) = 9$;
$\mathrm{Cov}(X,Y) = 1$, $\mathrm{Cov}(X,Z) = -1$, $\mathrm{Cov}(Y,Z) = 2$.
$\mathrm{Var}(X + Y + Z)$'i hesaplayın.
\end{alistirma}
\alcozum{%
$\mathrm{Var}(X+Y+Z) = \mathrm{Var}(X) + \mathrm{Var}(Y) + \mathrm{Var}(Z) + 2\mathrm{Cov}(X,Y) + 2\mathrm{Cov}(X,Z) + 2\mathrm{Cov}(Y,Z)$
$= 1 + 4 + 9 + 2(1) + 2(-1) + 2(2) = 14 + 2 - 2 + 4 = 18$.}

\begin{alistirma}[Al.4.9]{B}{Karakteristik Fonksiyon}
$X \sim \mathrm{Bernoulli}(p)$.
(a) $\varphi_X(t) = 1 - p + pe^{it}$ olduğunu gösterin.
(b) $X_1, \ldots, X_n$ bağımsız $\mathrm{Bernoulli}(p)$; $S_n = \sum X_i$.
$\varphi_{S_n}(t)$'yi bulun ve bunun hangi dağılıma ait olduğunu tanıyın.
\end{alistirma}
\alcozum{%
(a) $\varphi_X(t) = E[e^{itX}] = (1-p)e^{0} + pe^{it} = 1-p+pe^{it}$.\\
(b) Bağımsız çarpım: $\varphi_{S_n}(t) = (1-p+pe^{it})^n = \varphi_{\mathrm{Binom}(n,p)}(t)$.
Bu binom dağılımının KF'sidir; $S_n \sim \Binom(n,p)$.}

\begin{alistirma}[Al.4.10]{B}{Kule Özelliği}
$N \sim \mathrm{Pois}(\lambda)$ ve $X_1, X_2, \ldots$ bağımsız, $E[X_i] = \mu$, $\mathrm{Var}(X_i) = \sigma^2$.
$S_N = \sum_{i=1}^N X_i$ (toplam öğe sayısı rastlantısal).
(a) $E[S_N]$'i bulun. (b) $\mathrm{Var}(S_N)$'i bulun.
\end{alistirma}
\alcozum{%
(a) $E[S_N] = E[E[S_N \mid N]] = E[N\mu] = \mu E[N] = \mu\lambda$.\\
(b) Varyans ayrışım formülü:
$\mathrm{Var}(S_N) = E[\mathrm{Var}(S_N \mid N)] + \mathrm{Var}(E[S_N \mid N])
= E[N\sigma^2] + \mathrm{Var}(N\mu) = \sigma^2\lambda + \mu^2\lambda = \lambda(\sigma^2+\mu^2)$.}

\begin{alistirma}[Al.4.11]{C}{İspat: Cauchy-Schwarz Genel Formu}
$X, Y \in L^2(\Omega, \mathcal{F}, P)$ için $(E[XY])^2 \leq E[X^2] E[Y^2]$ olduğunu ispatlayın.
Eşitlik ne zaman sağlanır?
\end{alistirma}
\alcozum{%
Her $t \in \mathbb{R}$ için $f(t) = E[(tX+Y)^2] = t^2E[X^2] + 2tE[XY] + E[Y^2] \geq 0$.
Bu kuadratik $t$'de negatif olamaz; diskriminant $\leq 0$:
$4(E[XY])^2 - 4E[X^2]E[Y^2] \leq 0$.
Eşitlik: $f(t^*) = 0$ yani $t^*X + Y = 0$ n.k., yani $Y = cX$ n.k.\ (lineer bağımlılık).}

\begin{alistirma}[Al.4.12]{C}{Varyans Ayrışım Formülü (Kule Varyansı)}
$E[X^2] < \infty$ ve $\mathcal{G} \subseteq \mathcal{F}$ alt $\sigma$-cebiri olsun.
$\mathrm{Var}(X) = E[\mathrm{Var}(X \mid \mathcal{G})] + \mathrm{Var}(E[X \mid \mathcal{G}])$
olduğunu ispatlayın.
\end{alistirma}
\alcozum{%
$\mathrm{Var}(X) = E[X^2] - (E[X])^2$.
$E[\mathrm{Var}(X\mid\mathcal{G})] = E[E[X^2\mid\mathcal{G}] - (E[X\mid\mathcal{G}])^2]
= E[X^2] - E[(E[X\mid\mathcal{G}])^2]$.
$\mathrm{Var}(E[X\mid\mathcal{G}]) = E[(E[X\mid\mathcal{G}])^2] - (E[E[X\mid\mathcal{G}]])^2
= E[(E[X\mid\mathcal{G}])^2] - (E[X])^2$.
Toplam: $E[X^2] - (E[X])^2 = \mathrm{Var}(X)$. \checkmark}

\begin{alistirma}[Al.4.13]{A}{Chebyshev ve Cantelli Karşılaştırması}
$X$ rasgele değişken, $E[X] = 2$, $\mathrm{Var}(X) = 1$.
(a) Chebyshev ile $P(|X - 2| \geq 2)$ için üst sınır verin.
(b) Cantelli ile $P(X - 2 \geq 2)$ için üst sınır verin.
(c) Sonuçları karşılaştırın.
\end{alistirma}
\alcozum{%
(a) Chebyshev: $P(|X-2|\geq2) \leq \mathrm{Var}(X)/4 = 1/4$.\\
(b) Cantelli: $P(X-2\geq2) \leq \sigma^2/(\sigma^2+\lambda^2) = 1/(1+4) = 1/5$.\\
(c) Cantelli tek taraflı sınır veriyor ve daha keskin: $1/5 < 1/4$.
Chebyshev iki taraflı; tam $P(X-2\geq2)$ için Cantelli daha iyi.}

\begin{alistirma}[Al.4.14]{B}{Hoeffding Uygulaması}
$X_1, \ldots, X_{100}$ bağımsız, $X_i \in [0,1]$, $E[X_i] = 0.5$.
(a) Hoeffding eşitsizliği ile $P(\bar{X}_{100} \geq 0.6)$ için üst sınır verin.
(b) Chebyshev ile (varyanslı $\sigma^2 = 1/12$ alarak) karşılaştırın.
(c) $P \leq 0.01$ için gereken minimum $n$'yi Hoeffding ile bulun ($\varepsilon = 0.1$).
\end{alistirma}
\alcozum{%
(a) Hoeffding: $P(\bar{X}_n - 0.5 \geq 0.1) \leq e^{-2n(0.1)^2/1} = e^{-2(100)(0.01)} = e^{-2} \approx 0.135$.\\
(b) Chebyshev: $P(|\bar{X}_{100}-0.5|\geq0.1) \leq \sigma^2/(n\varepsilon^2) = (1/12)/(100\cdot0.01) = 1/12 \approx 0.083$.
Bu örnekte Chebyshev daha iyi — Hoeffding (sınırlı) avantajı daha büyük $n$'de belirginleşir.\\
(c) $e^{-2n\varepsilon^2} \leq 0.01$: $-2n(0.01) \leq \ln(0.01) = -4.605$; $n \geq 230$.}

\begin{alistirma}[Al.4.15]{B}{Çarpıklık ve Basıklık}
$X \sim \Exp(1)$: $E[X] = 1$, $\mathrm{Var}(X) = 1$.
(a) $E[X^3]$ ve $E[X^4]$'ü hesaplayın ($E[X^k] = k!$ Üstel için).
(b) Çarpıklık $\gamma_1 = (E[(X-\mu)^3])/\sigma^3$'ü hesaplayın.
(c) Basıklık fazlasını $\gamma_2 = E[(X-\mu)^4]/\sigma^4 - 3$ hesaplayın.
\end{alistirma}
\alcozum{%
(a) $E[X^3] = 3! = 6$; $E[X^4] = 4! = 24$.\\
(b) $\mu_3' = E[(X-1)^3] = E[X^3 - 3X^2 + 3X - 1] = 6 - 3(2) + 3(1) - 1 = 6-6+3-1 = 2$.
$\gamma_1 = 2/1^3 = 2$: kuvvetli sağ çarpıklık.\\
(c) $\mu_4' = E[(X-1)^4] = E[X^4 - 4X^3 + 6X^2 - 4X + 1] = 24 - 24 + 12 - 4 + 1 = 9$.
$\gamma_2 = 9/1 - 3 = 6$: kuvvetli ağır kuyruk (Normal'den 6 kat fazla).}

% =====================================================================
\endinput
